\documentclass[letterpaper,openright,12pt]{report}
\usepackage[spanish]{babel} % espanol
\usepackage[utf8]{inputenc} % acentos sin codigo
\usepackage{graphicx} % graficos
\usepackage{picture}

\begin{document}

\begin{titlepage}
\setlength{\unitlength}{1 cm} %Especificar unidad de trabajo
\begin{center}
\vspace*{-1in}
\begin{picture}(18,4)
%\begin{center}
\vspace*{0.18in}
\put(0,0){\includegraphics[width=3cm,height=4cm]{./.imagen/Logoutem.jpg}}
\put(11.5,0){\includegraphics[width=4cm,height=4cm]{./.imagen/iccilogocolor.png}}
%\includegraphics[width=4cm]{./.imagen/Logoutem.jpg}
%\end{center}
\end{picture}
\\[2cm]
\textbf{{\Huge Universidad Tecnológica Metropolitana}\\[1.5cm]
{\LARGE Escuela De informática}}\\[1.25cm]
{\LARGE \textbf{Busqueda Mediante Hashing}}\\[2.5cm]
{\large Integrantes:}\\
David Martínez Rivas\\
Felipe Canales Saavedra\\
Javier Reyes Gonzalez\\
Fernando Rubilar Zepeda\\
Profesor:\\
Alejandro Reyes\\[2cm]
Chile - \today
\end{center}

\end{titlepage}
\tableofcontents
\appendix

\newpage
\begin{center}
\section{Introducción}

En informática, hash se refiere a una función o método para generar
claves o llaves que representan de manera casi unívoca a un
documento,registr, arichivo, etc., resumir o identificar un dato a
tráves de la probabilidad, utilizando una función hash o algoritmo hash.
Un has es el resultado de dicha función o algoritmo
\end{center}
\newpage
\section{Historia}

\textbf{Orígenes del término} El término hash proviene, aparentemente,
de la analogía con el significado estándar (en inglés) de dicha palabra
en el mundo real: picar y mezclar. Donald Knuth cree que H. P. Luhn,
empleado de IBM, fue el primero en utilizar el concepto en un memorándum
fechado en enero de 1953. Su utilización masiva no fue hasta después de
10 años. En el algoritmo SHA-1, por ejemplo, el conjunto de partida de
la función es dividido en palabras que son mezcladas entre sí utilizando
funciones matemáticas seleccionadas especialmente. Se hace que el rango
de valores que puede devolver la función sea de longitud fija: 160 bits
utilizando la adición modular.
\newpage
\section{Función hash}

Es una función para resumir o identificar probabilísticamente un gran
conjunto de información, dando como resultado un conjunto imagen finito
generalmente menor. Varían en los conjuntos de partida y de llegada y en
cómo afectan a la salida similitudes o patrones de la entrada

\subsection{Ventajas}

Se pueden usar los valores naturales de la llave, puesto que se traducen
internamente a direcciones fáciles de localizar Se logra independencia
lógica y física, debido a que los valores de las llaves son
independientes del espacio de direcciones No se requiere almacenamiento
adicional para los índices.

\subsection{Desventajas}

El archivo no esta clasificado No permite llaves repetidas Solo permite
acceso por una sola llave Costos Tiempo de procesamiento requerido para
la aplicación de la función hash
\newpage

\section{Aplicación}

\begin{itemize}
\itemsep1pt\parskip0pt\parsep0pt
\item
  Es posible que existan claves resultantes iguales para objetos
  diferentes, ya que el rango de posibles claves es mucho menor que el
  de posibles objetos a resumir (las claves suelen tener en torno al
  centenar de bits, pero los ficheros no tienen un tamaño límite).
\item
  Son usadas en múltiples aplicaciones, como los arrays asociativos,
  criptografía, procesamiento de datos y firmas digitales, entre otros.
  Una buena función de hash es una que experimenta pocas colisiones en
  el conjunto esperado de entrada; es decir que se podrán identificar
  unívocamente las entradas.
\item
  Las tablas hash, una de las aplicaciones más extendidas de las
  funciones de hash, aceleran el proceso de búsqueda de un registro de
  información según una clave (nota: este uso de la palabra poco se
  relaciona con su significado habitual). Por ejemplo, una cadena
  alfanumérica puede ser utilizada para buscar la información de un
  empleado en la base de datos de un sistema.
\item
  La utilización de tablas hash provee un acceso casi directo a dichos
  registros, lo que significa que, en promedio, una búsqueda puede
  llegar a requerir sólo uno o dos intentos en la memoria o el archivo
  que contiene la información. Naturalmente, se prefiere una buena
  función de hash que evitará colisiones de hash.
\item
  Si asumimos que la clave es una cadena de bytes, entonces la función
  de hash debería ser como un índice de los registros que tiene una
  distribución aleatoria sobre las cadenas de entrada esperadas. De otra
  forma, habría más colisiones de hash degradando así el tiempo de
  búsqueda. Si, por ejemplo, la clave es alfabética, cada byte puede
  tener sólo 26 de sus 256 valores posibles. Funciones tan simples no
  distribuirán los índices de una forma pareja.
\item
  Para una comparación de velocidad y calidad de varias funciones de
  hash, referirse a los enlaces externos.
\end{itemize}
\newpage
\section{Conclusión}

El método de búsqueda hash o por transformación de clave aumenta la
velocidad de búsqueda sin necesidad de que los elementos estén
previamente ordenados.
\end{document}
